% =====================================================================
% main.tex --- EMNLP 2026 submission
% "A Scaling Law for Grammatical Acquisition: Dose-Response Curves for
%  Construction Learning in Small Language Models"
%
% --------------------------------------------------------------------
% STYLE FILES
% --------------------------------------------------------------------
% This paper targets the official ACL/EMNLP LaTeX style. Before the
% camera-ready (or for a faithful preview), download the EMNLP 2026
% release of:
%     acl.sty           (the document class wrapper / package)
%     acl_natbib.bst    (the bibliography style)
% from the ACL submission instructions page and drop them in THIS
% directory. With those present, \usepackage{acl} just works.
%
% If acl.sty is absent, the guarded block below falls back to a plain
% 11pt article preamble so the skeleton still compiles end-to-end. The
% fallback is for drafting only --- never submit it.
% --------------------------------------------------------------------
%
% NOTE ON TITLE: the working title below is neutral. The FINAL title is
% selected from the PRD Outcome Matrix once the pre-registered decision
% rule fires (see src/drc/analysis/decision.py and Appendix E). Do not
% lock the title before the decision script has run on real results.
% =====================================================================

\IfFileExists{acl.sty}{%
  \documentclass[11pt]{article}
  \usepackage[review]{acl}       % drop [review] for camera-ready
}{%
  % ---- Graceful fallback preamble (drafting only) ----
  \documentclass[11pt]{article}
  \usepackage[margin=1in]{geometry}
  \usepackage{natbib}
  \bibliographystyle{plainnat}
  \setlength{\parindent}{1em}
  \PackageWarningNoLine{main}{acl.sty not found --- using plain article
    fallback. Install the EMNLP 2026 style before submitting.}
}

\usepackage[T1]{fontenc}
\usepackage[utf8]{inputenc}
\usepackage{microtype}
\usepackage{graphicx}
\usepackage{booktabs}
\usepackage{amsmath}
\usepackage{amssymb}
\usepackage{xcolor}
\usepackage{url}

% --------------------------------------------------------------------
% PLACEHOLDER MACROS
% --------------------------------------------------------------------
% \resultpending{...} : a visible, impossible-to-miss stand-in for a
%   number or finding that depends on experiments NOT YET RUN. Every
%   data-dependent claim in this skeleton uses it. Search the compiled
%   PDF for "RESULT PENDING" to find everything still to fill.
\newcommand{\resultpending}[1]{%
  \textcolor{red}{\textbf{[RESULT PENDING: #1]}}}
% \TODO{...} : an author-facing note, also rendered visibly.
\newcommand{\TODO}[1]{\textcolor{blue}{\textbf{[TODO: #1]}}}

% Convenience: typeset the Hill equation parameters consistently.
\newcommand{\Ezero}{$E_0$}
\newcommand{\Emax}{$E_{\max}$}
\newcommand{\Efifty}{$\mathrm{EC}_{50}$}
\newcommand{\hilln}{$n$}

\title{A Scaling Law for Grammatical Acquisition:\\
  Dose-Response Curves for Construction Learning in\\
  Small Language Models}
% Working title --- final title chosen from the Outcome Matrix once the
% decision rule fires on real data (Appendix E).

\author{Anonymous Submission \\
  EMNLP 2026 \\
  \texttt{anonymous@example.org}}

\begin{document}
\maketitle

% =====================================================================
\begin{abstract}
How much exposure does a language model need to learn a grammatical
construction? We treat that question the way pharmacology treats a drug.
We hold a 10-million-word pretraining corpus fixed, then vary one thing:
the number of times a target construction appears, dosed at 0, 4, 16, 64,
and every attested instance. Train an LTG-BERT model from scratch at each
dose, measure acceptability with a syntactic log-odds-ratio (SLOR) score
on minimal pairs, and you can fit a dose-response curve and read off
interpretable parameters --- a floor with zero exposure, a ceiling, a
half-saturation dose, and a steepness that tells smooth scaling apart from
a sharp threshold. We run this across four constructions (AANN, the
comparative correlative, tough-movement, and the resultative) and three
random seeds, giving 60 models in a single frozen training recipe.
We pre-registered five hypotheses and the rules that pick between them.
\resultpending{which hypothesis fired and its one-line summary}
\resultpending{headline dose-response finding: mean Hill coefficient and
  what it implies about smooth vs.\ threshold learning}
\resultpending{cross-construction summary: do constructions cluster into
  distinct learning regimes?}
Our code, frozen corpora, and evaluation items are released for full
reproduction.
\end{abstract}

% =====================================================================
\section{Introduction}
\label{sec:intro}

A child hears some constructions thousands of times and others almost
never, yet ends up fluent in both. Language models face the same
lopsided diet. Some patterns saturate their training data; others, like
the article-adjective-numeral-noun construction in ``a beautiful five
days,'' show up only a handful of times in millions of words
\citep{mahowald2023aann,misra2024aann}. So a natural question follows:
how does a model's command of a construction depend on how often it sees
that construction during pretraining?

Most work answers this only at the extremes. Filtered-corpus training
asks whether a model can learn a phenomenon from \emph{zero} direct
exposure, leaning entirely on indirect evidence
\citep{patil2024fict,leong2024passives,misra2024aann}. That's a sharp,
useful test, but it skips the middle of the curve. The closest prior work,
\citet{oba2024widet}, does sample several exposure counts and is the right
precedent here --- yet it reads the resulting curves qualitatively
(``rises gradually, then plateaus'') and never fits a model or extracts a
number. That is exactly the gap we close. Between ``never seen it'' and
``seen it everywhere'' lies the interesting part: does competence rise
smoothly with exposure, or snap into place past a threshold? Does every
construction follow the same trajectory, or do different constructions
need wildly different amounts of data? We answer those questions with a
\emph{parametric} model whose fitted parameters are directly comparable
across constructions.

We borrow the workhorse tool of dose-response pharmacology. Swap drug
concentration for ``number of construction exposures in pretraining'' and
biological response for ``minimal-pair acceptability'' and the
four-parameter Hill equation maps cleanly onto our question. Its
parameters are readable. The floor \Ezero{} is what the model scores with
zero direct exposure --- a clean measure of how far indirect evidence
alone gets you. The ceiling \Emax{} is where competence saturates. The
half-saturation dose \Efifty{} answers ``how much data does this
construction need.'' And the Hill coefficient \hilln{} is the shape knob:
near 1 the curve is a gentle saturating rise, and much larger than that it
sharpens into a step --- the signature of a phase transition.

To make the curve mean something, we have to keep everything else fixed.
We start from BabyLM-10M \citep{warstadt2023babylm,choshen2024babylm} and,
for each construction, build a family of corpora that differ in exactly
one variable: the dose. Removed sentences are replaced by length- and
domain-matched sentences from a 100M-word pool, so corpus size and genre
mix stay constant while only the construction count moves. We then train
LTG-BERT \citep{samuel2023ltgbert} from scratch at each dose with one
frozen recipe, holding architecture, optimizer, and schedule identical
across all 60 runs. Tuning per dose would just invite worries about
p-hacking; freezing the recipe is what keeps the comparison honest.

We study four constructions chosen to span the space: AANN (productive but
rare), the comparative correlative (semi-idiomatic;
\citealp{weissweiler2022comparative}), tough-movement (a classic
poverty-of-stimulus case; \citealp{hu2020syntaxgym}), and the resultative
(an argument-structure construction; \citealp{goldberg2004resultatives}).
Three random seeds per cell let seed-to-seed spread become a real
distribution rather than a footnote.

Crucially, we locked our analysis before seeing any fitted numbers. Five
hypotheses --- smooth scaling, phase transition, heterogeneous regimes,
indirect-evidence dominance, and stochastic acquisition --- each map to a
specific paper framing, and a pre-registered decision rule picks the
winner from the data (Appendix~\ref{app:hypotheses}). The result we report
is whatever the rule selects, not a story we backfilled.

Our contributions:
\begin{itemize}
  \item The first \emph{parametric} dose-response characterization of
        construction learning: we fit Hill curves to acceptability as a
        function of pretraining exposure and extract a per-construction
        learnability threshold (\Efifty), steepness (\hilln), and an
        indirect-evidence floor (\Ezero) that are comparable across
        constructions --- turning the qualitative exposure curves of prior
        work \citep{oba2024widet} into quantitative estimates, on a frozen
        60-run design.
  \item Three analyses the dosing design makes possible and that prior work
        has not run: predicting \Efifty{} from corpus properties
        (Section~\ref{sec:predictability}), a cross-construction
        \emph{transfer} matrix (Section~\ref{sec:transfer}), and a
        seen-versus-novel-filler split that separates memorization from
        generalization (Section~\ref{sec:generalization}).
  \item A pre-registered hypothesis space and decision rule that fixes the
        interpretation in advance.
  \item Released code, frozen dose corpora, evaluation items, and fitted
        curves for full reproduction.
\end{itemize}

% =====================================================================
\section{Related Work}
\label{sec:related}

\paragraph{Filtered-corpus training and indirect evidence.}
The closest line of work removes a target phenomenon from pretraining and
asks whether a model learns it anyway. \citet{patil2024fict} formalize
this as filtered-corpus training and show models generalize to held-out
syntactic phenomena from indirect evidence alone. \citet{leong2024passives}
study the passive and find that learning from indirect evidence demands
large amounts of data. \citet{misra2024aann} run the cleanest case for our
purposes: they filter AANNs from training and show models still rate them,
because related, less-rare constructions teach the pattern indirectly.
Recent benchmark work probes the sample efficiency of these syntactic
heuristics directly \citep{anonymous2026posh}. Most of these fix the dose at
either zero or ``all''; the exception, \citet{oba2024widet}, injects a
phenomenon at several counts and is the work we build on most directly. We
differ in what we do with the resulting curve: rather than describing its
shape in prose, we fit a four-parameter model and extract comparable
learnability parameters, then use the dosing design to ask new questions
about transfer, predictability, and memorization.

\paragraph{The BabyLM setting.}
Training on a developmentally plausible budget of words has become its own
research program \citep{warstadt2023babylm,choshen2024babylm,rozner2025babylm}.
LTG-BERT \citep{samuel2023ltgbert} is a strong sample-efficient encoder in
this regime and is the architecture we adopt. The small-data setting is
ideal for dosing: a construction that appears a handful of times in 10M
words is genuinely rare, so the low-dose conditions are realistic rather
than artificial. Work on impossible languages \citep{kallini2024mission}
and on formal-language pre-pretraining \citep{hu2025circuits} shows these
small models are sensitive probes of structural learnability.

\paragraph{Constructions and acquisition.}
We take a construction-grammar view: form-meaning pairings that range from
fully productive to frozen and idiomatic
\citep{goldberg1995constructions,goldberg2004resultatives,tomasello2003constructing}.
This contrasts with a poverty-of-stimulus framing on which core grammar is
underdetermined by the input \citep{chomsky1965aspects}, a debate that
language models have reopened \citep{lan2026llmsps}. The four targets here
have established linguistic analyses and existing probes
\citep{weissweiler2022comparative,hu2020syntaxgym,goldberg2004resultatives,mahowald2023aann}.
Work comparing model and child learning stages motivates treating exposure
as a developmental variable \citep{evanson2023acquisition}.

\paragraph{Scaling, frequency, and memorization.}
Frequency drives what models learn: rare, long-tail knowledge needs far
more exposure \citep{kandpal2023longtail}, and recent work separates
generalization from memorization as a function of pretraining frequency
\citep{wang2024memorization,morris2025memorization}. We measure that
dependence directly for grammatical constructions and fit it with an
interpretable curve.

\paragraph{Learning curves over training, not exposure.}
A related strand fits curves to syntactic competence as a function of
\emph{training steps} on a fixed corpus: \citet{bunzeck2024shapes} find
non-uniform, sometimes non-monotone BLiMP learning curves, and
\citet{chang2024curves} fit sigmoids and define a 50\%-surprisal ``age of
acquisition.'' Our independent variable is different in kind --- the
\emph{count} of a construction in the data, not wall-clock training --- so
our \Efifty{} is an exposure threshold, not a temporal one. The two axes
are complementary.

\paragraph{Acceptability measurement.}
We score acceptability with SLOR \citep{pauls2012slor,lau2017acceptability},
which corrects a sentence's log-probability for word rarity, evaluated on
minimal pairs in the BLiMP tradition \citep{warstadt2020blimp}. For our
masked encoder we use the pseudo-log-likelihood variant standard for masked
LMs.

% =====================================================================
\section{Methods}
\label{sec:methods}

Everything in this section is fixed by design and does not depend on
results. The recipe is frozen across all 60 runs.

\subsection{Construction selection}
\label{sec:methods-constructions}
We study four English constructions chosen to span productivity,
idiomaticity, and argument structure:
\begin{itemize}
  \item \textbf{AANN}: article + adjective + numeral + plural noun, as in
        ``a beautiful five days in Texas.'' Productive but rare; the
        indefinite article scopes over a plural measured quantity
        \citep{mahowald2023aann,misra2024aann}.
  \item \textbf{Comparative correlative}: ``the X-er $\ldots$, the
        Y-er $\ldots$,'' as in ``the harder you try, the worse it gets.''
        Semi-idiomatic, with a frozen degree-marking ``the''
        \citep{weissweiler2022comparative}.
  \item \textbf{Tough-movement}: ``this book is tough to read,'' where the
        understood object of the embedded verb surfaces as the matrix
        subject. A textbook poverty-of-stimulus dependency
        \citep{hu2020syntaxgym}.
  \item \textbf{Resultative}: V + NP + AP where the AP names a result
        state, as in ``he hammered the metal flat.'' An argument-structure
        construction \citep{goldberg2004resultatives}.
\end{itemize}
Each construction is detected by a dedicated structural filter over
dependency parses, with surface fallbacks where parses are brittle; full
filter specifications are in Appendix~\ref{app:filters}.

\subsection{Dose corpus construction}
\label{sec:methods-dose}
For each construction we build five corpora that differ in exactly one
variable --- the number of attested instances of the target construction
present: doses 0, 4, 16, 64, and ``all'' (every instance left in place).
The challenge is holding everything else constant. We do not simply delete
positive sentences, since that would shrink the corpus and skew its genre
mix, confounding the dose with corpus size. Instead, every removed
sentence is replaced by a matched sentence drawn from a 100M-word pool,
selecting in priority order: (1) same source domain, (2) length within
$\pm 20\%$, relaxing to $\pm 40\%$, then (3) the nearest-length sentence
from the same domain. Corpus construction uses one fixed seed (42); the
dose corpora are generated once and frozen, and only training seeds vary
the model. Each corpus ships with an audit JSON recording every kept,
removed, and swapped-in sentence, and we re-run the construction filter
afterward to verify the realized counts. Parsing uses Stanza
\citep{qi2020stanza}.

\subsection{Models and training}
\label{sec:methods-models}
We train LTG-BERT \citep{samuel2023ltgbert} from scratch on each dose
corpus, using the BabyLM strict-track winning hyperparameters. The recipe
is frozen across all runs: 12 layers, hidden size 384, 6 attention heads,
intermediate size 1024, maximum sequence length 128, and a vocabulary of
16{,}384 (a WordPiece tokenizer trained per dose corpus). We optimize the
masked-LM objective with AdamW ($\beta_1{=}0.9$, $\beta_2{=}0.98$,
$\epsilon{=}10^{-6}$, weight decay 0.1), peak learning rate $10^{-3}$ with
a 10\% linear warmup and cosine decay, batch size 256, 20 epochs, bf16
precision, and 15\% whole-word masking. We carve a 100k-word held-out
slice off each corpus to track perplexity. The full hyperparameter table
is in Appendix~\ref{app:hyperparams}. We run three random seeds (42, 43,
44) per (construction, dose) cell, for $4 \times 5 \times 3 = 60$ models.
Implementation uses Hugging Face Transformers \citep{wolf2020transformers}.

\subsection{SLOR evaluation}
\label{sec:methods-slor}
We measure acceptability with the syntactic log-odds-ratio
\citep{pauls2012slor,lau2017acceptability}, which corrects a sentence's
log-probability for the fact that rare words drag it down regardless of
grammaticality:
\begin{equation}
  \mathrm{SLOR}(s) = \frac{\log P(s) - \log P_{\mathrm{uni}}(s)}{|s|},
\end{equation}
where $P_{\mathrm{uni}}$ is a corpus unigram model and $|s|$ is the word
count. For our masked encoder there is no left-to-right $\log P(s)$, so we
use the pseudo-log-likelihood: mask each content position in turn, read off
the log-probability of the true token given the rest, and sum. The unigram
correction for each model is built from that model's \emph{own} dose
corpus, so the normalizer matches the word-frequency world the model
trained on. On a minimal pair, the model is correct when
$\mathrm{SLOR}(\text{good}) > \mathrm{SLOR}(\text{bad})$. We report
per-(model, construction) mean accuracy with a binomial standard error.
Each trained model is scored against \emph{all four} construction test
sets, not just its own, so cross-construction (collateral) effects surface
rather than hide. Evaluation item schemas and counts are in
Appendix~\ref{app:items}. We also fit an $n$-gram baseline
\citep{warstadt2020blimp} as a control.

\subsection{Statistical analysis and Hill fitting}
\label{sec:methods-hill}
For each (construction, seed) we fit the four-parameter Hill equation to
self-evaluation accuracy as a function of dose:
\begin{equation}
  Y(D) = E_0 + (E_{\max} - E_0)\,\frac{D^{\,n}}{\mathrm{EC}_{50}^{\,n} + D^{\,n}}.
  \label{eq:hill}
\end{equation}
We fit with bounded nonlinear least squares (trust-region reflective),
constraining \Ezero{} and \Emax{} to $[0,1]$ and \Efifty{}, \hilln{} to be
positive. The dose ``all'' is resolved to the construction's attested count
recorded at generation time. Confidence intervals come from a parametric
bootstrap: resample residuals under the fitted curve, refit, and take 2.5
/97.5 percentiles over 1{,}000 resamples \citep{efron1993bootstrap}. The
bootstrap avoids leaning on asymptotic normality, which is shaky with only
five dose points.

A good Hill fit alone proves little, so we pit it against simpler shapes
and let AIC and BIC referee: a power law $a(D{+}1)^b{+}c$, a step function,
a log-linear $a + b\log(D{+}1)$, and a null (the mean). To ask whether
constructions learn the same way, we cluster the fitted parameter vectors
$[E_0, \log \mathrm{EC}_{50}, n, E_{\max}]$ with $k$-means for
$k \in \{2,3,4\}$ and record silhouette scores. A clean split at $k{=}2$ is
the quantitative tell for distinct learning regimes.

Finally, we apply the pre-registered decision rule
(Appendix~\ref{app:hypotheses}), which reads these statistics and names the
hypothesis the data supports. The rule, the thresholds, and the
hypothesis-to-title mapping were all locked before any fit was inspected.

% =====================================================================
\section{Results}
\label{sec:results}

% All quantitative content in this section is pending experiments. Tables
% and the decision verdict will be filled from results/*.csv and
% results/decision.txt; figures are auto-generated into results/figures/.

\subsection{Dose-response curves}
\label{sec:results-curves}
Figure~\ref{fig:dose-response} shows fitted dose-response curves per
construction, with per-seed lines, the mean, a bootstrap ribbon, and the
Hill overlay. \resultpending{describe the overall curve shape and whether
fits converged across cells; fill from \texttt{results/hill\_fits.csv}}

\begin{figure}[t]
  \centering
  \IfFileExists{../results/figures/fig1_dose_response_curves.pdf}{%
    \includegraphics[width=\columnwidth]{../results/figures/fig1_dose_response_curves.pdf}}{%
    \fbox{\parbox{0.9\columnwidth}{\centering\vspace{2em}
      \resultpending{fig1\_dose\_response\_curves.pdf not yet generated ---
        run \texttt{python -m drc.analysis.figures}}\vspace{2em}}}}
  \caption{Dose-response curves by construction (2$\times$2). Per-seed
    lines, mean, bootstrap ribbon, and fitted Hill overlay.}
  \label{fig:dose-response}
\end{figure}

\subsection{Fitted Hill parameters}
\label{sec:results-params}
Table~\ref{tab:hill} reports the fitted Hill parameters per construction,
averaged over seeds with bootstrap 95\% CIs. The Hill coefficient \hilln{}
is the key quantity: it separates smooth scaling from a phase transition.

\begin{table}[t]
  \centering
  \small
  \begin{tabular}{lcccc}
    \toprule
    Construction & \Ezero & \Emax & \Efifty & \hilln \\
    \midrule
    AANN          & \resultpending{} & \resultpending{} & \resultpending{} & \resultpending{} \\
    Comp.\ corr.  & \resultpending{} & \resultpending{} & \resultpending{} & \resultpending{} \\
    Tough-mov.    & \resultpending{} & \resultpending{} & \resultpending{} & \resultpending{} \\
    Resultative   & \resultpending{} & \resultpending{} & \resultpending{} & \resultpending{} \\
    \bottomrule
  \end{tabular}
  \caption{Fitted Hill parameters (mean over seeds, 95\% bootstrap CI).
    Fill from \texttt{results/hill\_fits.csv}.}
  \label{tab:hill}
\end{table}

Figure~\ref{fig:clusters} scatters the fitted parameters, colored by
construction and marked by seed.

\begin{figure}[t]
  \centering
  \IfFileExists{../results/figures/fig2_parameter_clusters.pdf}{%
    \includegraphics[width=\columnwidth]{../results/figures/fig2_parameter_clusters.pdf}}{%
    \fbox{\parbox{0.9\columnwidth}{\centering\vspace{2em}
      \resultpending{fig2\_parameter\_clusters.pdf not yet generated}\vspace{2em}}}}
  \caption{Fitted Hill parameters per construction and seed.}
  \label{fig:clusters}
\end{figure}

\subsection{Model comparison and clustering}
\label{sec:results-modelcomp}
Table~\ref{tab:modelcomp} reports, per construction, which curve shape wins
by AIC, testing whether the Hill curve earns its parameters against simpler
alternatives. \resultpending{which model wins per construction and the
$k{=}2$ silhouette; fill from results/model\_comparison.csv and
results/cluster\_assignments.csv}

\begin{table}[t]
  \centering
  \small
  \begin{tabular}{lcc}
    \toprule
    Construction & Best model (AIC) & $\Delta$AIC vs.\ next \\
    \midrule
    AANN          & \resultpending{} & \resultpending{} \\
    Comp.\ corr.  & \resultpending{} & \resultpending{} \\
    Tough-mov.    & \resultpending{} & \resultpending{} \\
    Resultative   & \resultpending{} & \resultpending{} \\
    \bottomrule
  \end{tabular}
  \caption{Curve-shape model comparison. Fill from
    \texttt{results/model\_comparison.csv}.}
  \label{tab:modelcomp}
\end{table}

\subsection{Seed variance}
\label{sec:results-seed}
Figure~\ref{fig:seedvar} shows a construction-by-dose heatmap of the
coefficient of variation across seeds, the input to the stochastic
hypothesis test. \resultpending{whether within-seed wobble swamps the dose
effect}

\begin{figure}[t]
  \centering
  \IfFileExists{../results/figures/fig3_seed_variance.pdf}{%
    \includegraphics[width=\columnwidth]{../results/figures/fig3_seed_variance.pdf}}{%
    \fbox{\parbox{0.9\columnwidth}{\centering\vspace{2em}
      \resultpending{fig3\_seed\_variance.pdf not yet generated}\vspace{2em}}}}
  \caption{Coefficient of variation across seeds, by construction and dose.}
  \label{fig:seedvar}
\end{figure}

\subsection{Baseline and collateral effects}
\label{sec:results-collateral}
Figure~\ref{fig:ngram} compares the neural LM against the $n$-gram control,
and Figure~\ref{fig:collateral} shows the 4$\times$4 train-by-eval accuracy
heatmap that exposes collateral effects of dosing one construction on
another. \resultpending{do off-diagonal effects appear?}

\begin{figure}[t]
  \centering
  \IfFileExists{../results/figures/fig4_ngram_comparison.pdf}{%
    \includegraphics[width=\columnwidth]{../results/figures/fig4_ngram_comparison.pdf}}{%
    \fbox{\parbox{0.9\columnwidth}{\centering\vspace{2em}
      \resultpending{fig4\_ngram\_comparison.pdf not yet generated}\vspace{2em}}}}
  \caption{Neural LM vs.\ $n$-gram baseline (2$\times$2).}
  \label{fig:ngram}
\end{figure}

\begin{figure}[t]
  \centering
  \IfFileExists{../results/figures/fig5_collateral_effects.pdf}{%
    \includegraphics[width=\columnwidth]{../results/figures/fig5_collateral_effects.pdf}}{%
    \fbox{\parbox{0.9\columnwidth}{\centering\vspace{2em}
      \resultpending{fig5\_collateral\_effects.pdf not yet generated}\vspace{2em}}}}
  \caption{Collateral effects: mean accuracy across train/eval pairs.}
  \label{fig:collateral}
\end{figure}

\subsection{Predicting the threshold from corpus properties}
\label{sec:predictability}
If \Efifty{} is a real learnability measure, prior corpus properties should
predict it. We regress each construction's fitted \Efifty{} (and \Ezero,
\hilln) on its attested frequency, slot-filler productivity, and surface
unigram predictability (Figure~\ref{fig:predict}). With only four
constructions this is exploratory, and we report it as such.
\resultpending{which property tracks \Efifty, with the correlation; fill
from results/predictability\_correlations.csv}

\begin{figure}[t]
  \centering
  \IfFileExists{../results/figures/fig6_predictability.pdf}{%
    \includegraphics[width=\columnwidth]{../results/figures/fig6_predictability.pdf}}{%
    \fbox{\parbox{0.9\columnwidth}{\centering\vspace{2em}
      \resultpending{fig6\_predictability.pdf not yet generated}\vspace{2em}}}}
  \caption{Fitted \Efifty{} against the strongest corpus predictor
    (constructions, mean over seeds). Exploratory, $n=4$.}
  \label{fig:predict}
\end{figure}

\subsection{Cross-construction transfer}
\label{sec:transfer}
Because every model is evaluated on all four test sets, we can ask whether
dosing one construction shifts acceptability on another.
Figure~\ref{fig:transfer} reports the slope of each evaluated
construction's accuracy against the (log) dose of each trained
construction; the diagonal is the on-target effect, the off-diagonal is
transfer. \resultpending{do off-diagonal slopes depart from zero?}

\begin{figure}[t]
  \centering
  \IfFileExists{../results/figures/fig7_transfer_matrix.pdf}{%
    \includegraphics[width=\columnwidth]{../results/figures/fig7_transfer_matrix.pdf}}{%
    \fbox{\parbox{0.9\columnwidth}{\centering\vspace{2em}
      \resultpending{fig7\_transfer\_matrix.pdf not yet generated}\vspace{2em}}}}
  \caption{Transfer matrix: slope of eval-construction accuracy vs.\ the
    dose of the trained construction.}
  \label{fig:transfer}
\end{figure}

\subsection{Memorization versus generalization}
\label{sec:generalization}
We split each construction's test items by whether their lexical material
overlaps the instances kept at a given dose, then fit dose-response
separately for ``seen'' and ``novel'' fillers (Figure~\ref{fig:gen}). A
novel-filler curve that still climbs with dose is evidence the model
generalizes the construction rather than memorizing specific items.
\resultpending{does the novel-filler curve track the seen-filler curve?}

\begin{figure}[t]
  \centering
  \IfFileExists{../results/figures/fig8_generalization.pdf}{%
    \includegraphics[width=\columnwidth]{../results/figures/fig8_generalization.pdf}}{%
    \fbox{\parbox{0.9\columnwidth}{\centering\vspace{2em}
      \resultpending{fig8\_generalization.pdf not yet generated}\vspace{2em}}}}
  \caption{Dose-response for seen vs.\ novel lexical fillers, per
    construction. The gap indexes memorization.}
  \label{fig:gen}
\end{figure}

\subsection{Pre-registered decision}
\label{sec:results-decision}
Applying the locked decision rule
(Appendix~\ref{app:hypotheses}) to the statistics above:
\resultpending{which of H1--H5 fired, with the supporting numbers; fill
from results/decision.txt}

% =====================================================================
\section{Discussion}
\label{sec:discussion}

% The framing below is keyed to which hypothesis fires. Keep only the
% matching block once results/decision.txt names the winner; the others
% are guarded scaffolding so no claim is made prematurely.

\resultpending{which hypothesis fired --- keep the matching paragraph, delete the rest}

\paragraph{If H1 (smooth scaling) fired.}
% Hill wins by AIC for every construction and mean n in [0.5, 2.0].
A single saturating law would describe construction learning: competence
rises gradually and predictably with exposure, and the half-saturation
dose \Efifty{} becomes a usable ``data budget'' per construction.
\resultpending{quantify and contrast \Efifty{} across the four constructions}

\paragraph{If H2 (phase transition) fired.}
% Hill wins everywhere but mean n > 4: grammatical grokking.
The same curve, far steeper, would point to grammatical grokking ---
competence snapping into place past a threshold rather than accumulating.
\resultpending{report the threshold dose and steepness}

\paragraph{If H3 (heterogeneous regimes) fired.}
% Clean k=2 split, or Hill wins for some and step for others.
Constructions would split into distinct learning regimes, not one law.
\resultpending{name which constructions fall in which regime and why}

\paragraph{If H4 (indirect evidence dominates) fired.}
% Mean (Emax - E0)/Emax < 0.1: direct exposure barely moves the needle.
Direct exposure would barely matter: a high floor \Ezero{} means indirect
evidence already does most of the work, echoing filtered-corpus findings
\citep{patil2024fict,misra2024aann}.
\resultpending{report the floor-to-ceiling gain per construction}

\paragraph{If H5 (stochastic acquisition) fired.}
% Within-seed CV exceeds between-dose CV.
Seed-to-seed variation would swamp the dose effect, making acquisition
look more like a coin flip on initialization than a function of exposure.
\resultpending{report within- vs.\ between-cell CV}

\paragraph{If the rule was mixed.}
\resultpending{summarize the mixed verdict and point to the Outcome Matrix}

% =====================================================================
\section{Conclusion}
\label{sec:conclusion}
We framed construction learning as a dose-response problem, fit
interpretable Hill curves to acceptability as a function of pretraining
exposure, and pre-registered the rule that interprets them. The design ---
a frozen 60-run sweep with matched-replacement dose corpora --- isolates
exposure from corpus size and genre. \resultpending{one-sentence takeaway
once the decision rule has fired}

% =====================================================================
\section*{Limitations}
\label{sec:limitations}
Our study is deliberately narrow. We train small encoders (about 24M
parameters) on a 10M-word budget; whether the same curves hold for larger
models and corpora is open. We study four English constructions, so claims
about ``construction learning'' generalize only as far as that sample. SLOR
on minimal pairs is one acceptability proxy among several, and the
pseudo-log-likelihood for masked LMs is an approximation. Doses top out at
the attested count in 10M words, so we cannot probe the very high-exposure
regime that large web corpora reach. Five dose points is few for nonlinear
curve fitting, which is exactly why we bootstrap rather than trust
asymptotic errors, and why we test the Hill curve against simpler shapes.
Matched replacement controls corpus size and domain but cannot guarantee
that no indirect cue to the target construction rides along in the
replacement sentences. Finally, three seeds give a coarse estimate of
seed variance.

% =====================================================================
\section*{Ethics Statement}
\label{sec:ethics}
This work trains small language models from scratch on the public
BabyLM-10M corpus and a public 100M-word replacement pool, using existing
released data; we introduce no new human-subjects data and no personally
identifying information. The compute footprint is modest by design: 60
small-model pretraining runs on a developmentally plausible word budget,
in keeping with the sample-efficiency goals of the BabyLM program. Our
released artifacts (code, frozen corpora, evaluation items, fitted curves)
are intended to support reproduction and carry the same low risk profile as
the underlying public data. We see no direct path to harm beyond the
general risks of language modeling.

% =====================================================================
\bibliography{references}

\end{document}
